% Shared preamble for all AI Engineering Lab pre/post survey documents.
% Documents using this must declare \documentclass[addpoints,11pt]{exam} themselves,
% and may set \printanswers before \input-ing this file.

\usepackage[utf8]{inputenc}
\usepackage[T1]{fontenc}
\usepackage[english]{babel}
\usepackage{lmodern}
\usepackage{amssymb}
\usepackage{array}
\usepackage{tabularx}
\usepackage{booktabs}
\usepackage{enumitem}
\usepackage{xcolor}
\usepackage{tikz}
\usepackage{needspace}
\usepackage[hidelinks]{hyperref}

\usepackage[a4paper,margin=2.2cm,headsep=0.6cm,footskip=1.1cm]{geometry}

% ---------------------------------------------------------------------------
% Document metadata. Each driver file sets these before % Shared preamble for all AI Engineering Lab pre/post survey documents.
% Documents using this must declare \documentclass[addpoints,11pt]{exam} themselves,
% and may set \printanswers before \input-ing this file.

\usepackage[utf8]{inputenc}
\usepackage[T1]{fontenc}
\usepackage[english]{babel}
\usepackage{lmodern}
\usepackage{amssymb}
\usepackage{array}
\usepackage{tabularx}
\usepackage{booktabs}
\usepackage{enumitem}
\usepackage{xcolor}
\usepackage{tikz}
\usepackage{needspace}
\usepackage[hidelinks]{hyperref}

\usepackage[a4paper,margin=2.2cm,headsep=0.6cm,footskip=1.1cm]{geometry}

% ---------------------------------------------------------------------------
% Document metadata. Each driver file sets these before % Shared preamble for all AI Engineering Lab pre/post survey documents.
% Documents using this must declare \documentclass[addpoints,11pt]{exam} themselves,
% and may set \printanswers before \input-ing this file.

\usepackage[utf8]{inputenc}
\usepackage[T1]{fontenc}
\usepackage[english]{babel}
\usepackage{lmodern}
\usepackage{amssymb}
\usepackage{array}
\usepackage{tabularx}
\usepackage{booktabs}
\usepackage{enumitem}
\usepackage{xcolor}
\usepackage{tikz}
\usepackage{needspace}
\usepackage[hidelinks]{hyperref}

\usepackage[a4paper,margin=2.2cm,headsep=0.6cm,footskip=1.1cm]{geometry}

% ---------------------------------------------------------------------------
% Document metadata. Each driver file sets these before % Shared preamble for all AI Engineering Lab pre/post survey documents.
% Documents using this must declare \documentclass[addpoints,11pt]{exam} themselves,
% and may set \printanswers before \input-ing this file.

\usepackage[utf8]{inputenc}
\usepackage[T1]{fontenc}
\usepackage[english]{babel}
\usepackage{lmodern}
\usepackage{amssymb}
\usepackage{array}
\usepackage{tabularx}
\usepackage{booktabs}
\usepackage{enumitem}
\usepackage{xcolor}
\usepackage{tikz}
\usepackage{needspace}
\usepackage[hidelinks]{hyperref}

\usepackage[a4paper,margin=2.2cm,headsep=0.6cm,footskip=1.1cm]{geometry}

% ---------------------------------------------------------------------------
% Document metadata. Each driver file sets these before \input{common/preamble}.
%   \surveywave  -> "Pre-Test" / "Post-Test"
%   \surveyform  -> "A" / "B"
% ---------------------------------------------------------------------------
\providecommand{\surveywave}{Survey}
\providecommand{\surveyform}{A}

\newcommand{\labtitle}{AI Engineering Lab}

% ---------------------------------------------------------------------------
% Running head / foot
% ---------------------------------------------------------------------------
\pagestyle{headandfoot}
\firstpageheader{}{}{}
\runningheader{\small\labtitle}{}{\small\surveywave{} --- Form \surveyform}
\runningheadrule
\firstpagefooter{}{\small Page \thepage\ of \numpages}{}
\runningfooter{}{\small Page \thepage\ of \numpages}{}

% No per-item point values: every quiz item is scored 1/0 externally
% (see docs/analysis_plan.md), so the printed form carries no point marks.
\nopointsinmargin

% ---------------------------------------------------------------------------
% Section headings for the questionnaire parts
% ---------------------------------------------------------------------------
\newcommand{\partheading}[2]{%
  \par\vspace{1.1em}%
  \noindent\begin{tabularx}{\linewidth}{@{}>{\bfseries}l X@{}}
    \large Part #1 & \large\bfseries #2
  \end{tabularx}%
  \par\vspace{-0.4em}\noindent\rule{\linewidth}{1pt}\par\vspace{0.5em}%
}

\newcommand{\blockheading}[1]{%
  \par\vspace{0.8em}\noindent\textbf{#1}%
  \par\vspace{-0.3em}\noindent\rule{\linewidth}{0.4pt}\par\vspace{0.3em}%
}

\newcommand{\instruction}[1]{\par\vspace{0.2em}\noindent\textit{\small #1}\par\vspace{0.5em}}

% A block heading *inside* the questions environment. Plain \blockheading cannot
% be used there: the questions environment is a list, and material placed in a
% list before or between items has to be escaped through \fullwidth.
\newcommand{\quizblock}[1]{\fullwidth{\blockheading{#1}}}

% ---------------------------------------------------------------------------
% Answer primitives
% ---------------------------------------------------------------------------
\newcommand{\bx}{$\square$}                        % empty checkbox
\newcommand{\bxl}[1]{\bx~#1}                       % checkbox with label
\newcommand{\writeline}[1][\linewidth]{\rule[-0.4em]{#1}{0.4pt}}

% Single-character write-in boxes, e.g. for the pseudonym code.
\newcommand{\charbox}{%
  \tikz[baseline=(c.base)]{\node[draw,minimum width=8mm,minimum height=9mm,inner sep=0pt] (c) {\phantom{X}};}%
}
\newcommand{\charboxes}[1]{%
  \foreach \i in {1,...,#1} {\charbox\hspace{1.2mm}}%
}

% ---------------------------------------------------------------------------
% Rating grids
% ---------------------------------------------------------------------------
% These use plain `tabular` with an explicitly computed p{} label column rather
% than `tabularx`. tabularx determines its body by scanning the input for a
% literal \end{tabularx}, which it cannot find when the \begin/\end pair is
% split across an environment definition -- so tabularx must not be wrapped in
% a \newenvironment. The label width is therefore derived from \linewidth minus
% the total width of the fixed check-box columns.

\setlength{\extrarowheight}{0pt}
\setlength{\emergencystretch}{2em}
\newlength{\gridlabel}
\newcommand{\ck}{\makebox[8mm]{$\square$}}          % one check-box column
\newcommand{\hd}[1]{\makebox[8mm]{\scriptsize #1}}  % matching header cell

% Forced-choice 1-5 rating scale. There is deliberately no "n/a" / "cannot judge"
% column: every statement here is about the respondent's own knowledge or
% experience, which they are always in a position to place on the scale. An
% escape column on such a grid is used mainly to skip, and each skip removes a
% repeated measure that cannot be recovered afterwards. A genuinely unanswerable
% row can still be left blank, which is recorded as missing.
%
% The legend is a mandatory argument because the grids do not all use the same
% scale: agreement and ability need different anchors, and a printed legend that
% does not match the question is worse than none.
% Usage: \begin{likert}[<column header>]{<legend>} \litem{stem} ... \end{likert}
\newcommand{\legendagree}{1 = strongly disagree \quad 2 = disagree \quad 3 = neutral \quad 4 = agree
  \quad 5 = strongly agree}
\newcommand{\legendability}{1 = no experience \quad 2 = beginner \quad 3 = competent
  \quad 4 = proficient \quad 5 = expert}
\newenvironment{likert}[2][Today, I would say \dots]{%
  \par\small
  % The legend is emitted after the table, where the environment's arguments are
  % no longer in scope, so it is stashed here.
  \def\likertlegend{#2}%
  \setlength{\tabcolsep}{2pt}%
  \setlength{\gridlabel}{\dimexpr\linewidth-56mm\relax}%
  \renewcommand{\arraystretch}{1.3}%
  \noindent
  \tabular{@{}p{\gridlabel} ccccc@{}}
    \toprule
    \textbf{#1} & \hd{1} & \hd{2} & \hd{3} & \hd{4} & \hd{5} \\
    \midrule
}{%
    \bottomrule
  \endtabular
  \par\vspace{0.35em}%
  {\scriptsize \likertlegend}%
  \par
}
\newcommand{\litem}[1]{#1 & \ck & \ck & \ck & \ck & \ck \\}

% Two-scale component grid: "can explain" and "can implement" side by side.
\newenvironment{dualgrid}[2]{%
  \par\small
  \setlength{\tabcolsep}{1.6pt}%
  \setlength{\gridlabel}{\dimexpr\linewidth-97mm\relax}%
  \renewcommand{\arraystretch}{1.3}%
  \noindent
  \tabular{@{}p{\gridlabel} @{\hspace{2mm}} ccccc @{\hspace{4mm}} ccccc@{}}
    \toprule
    & \multicolumn{5}{c}{\textbf{#1}} & \multicolumn{5}{c}{\textbf{#2}} \\
    \cmidrule(lr){2-6}\cmidrule(lr){7-11}
    \textbf{Component} & \hd{1} & \hd{2} & \hd{3} & \hd{4} & \hd{5}
                       & \hd{1} & \hd{2} & \hd{3} & \hd{4} & \hd{5} \\
    \midrule
}{%
    \bottomrule
  \endtabular
  \par\vspace{0.35em}%
  {\scriptsize 1 = not at all \quad 2 = a little \quad 3 = partly \quad 4 = mostly \quad 5 = fully}%
  \par
}
\newcommand{\ditem}[1]{#1 & \ck & \ck & \ck & \ck & \ck & \ck & \ck & \ck & \ck & \ck \\}

% Frequency grid: daily / weekly / monthly / tried once / never.
\newenvironment{freqgrid}[1]{%
  \par\small
  \setlength{\tabcolsep}{2pt}%
  \setlength{\gridlabel}{\dimexpr\linewidth-75mm\relax}%
  \renewcommand{\arraystretch}{1.3}%
  \noindent
  \tabular{@{}p{\gridlabel} ccccc@{}}
    \toprule
    \textbf{#1} & \makebox[12mm]{\scriptsize daily} & \makebox[12mm]{\scriptsize weekly}
                & \makebox[14mm]{\scriptsize monthly} & \makebox[16mm]{\scriptsize tried once}
                & \makebox[12mm]{\scriptsize never} \\
    \midrule
}{%
    \bottomrule
  \endtabular\par
}
\newcommand{\fitem}[1]{#1 & \makebox[12mm]{\bx} & \makebox[12mm]{\bx} & \makebox[14mm]{\bx}
                          & \makebox[16mm]{\bx} & \makebox[12mm]{\bx} \\}

% Experience grid: yes / attempted / no.
\newenvironment{expgrid}[1]{%
  \par\small
  \setlength{\tabcolsep}{2pt}%
  \setlength{\gridlabel}{\dimexpr\linewidth-54mm\relax}%
  \renewcommand{\arraystretch}{1.3}%
  \noindent
  \tabular{@{}p{\gridlabel} ccc@{}}
    \toprule
    \textbf{#1} & \makebox[12mm]{\scriptsize yes} & \makebox[22mm]{\scriptsize attempted}
                & \makebox[12mm]{\scriptsize no} \\
    \midrule
}{%
    \bottomrule
  \endtabular\par
}
\newcommand{\eitem}[1]{#1 & \makebox[12mm]{\bx} & \makebox[22mm]{\bx} & \makebox[12mm]{\bx} \\}

% ---------------------------------------------------------------------------
% Quiz items
% ---------------------------------------------------------------------------
% The "don't know" escape hatch. Placed after the choices so that the item
% itself remains a four-option single-best-answer item, while omission is
% still distinguishable from guessing in the data.
\newcommand{\dontknow}{%
  \par\vspace{-0.2em}\noindent\hspace{1.2em}{\small\bx~\textit{I do not know.}}\par
}

% Metadata shown only in the answer key.
\newcommand{\itemmeta}[4]{%
  \ifprintanswers
    \par\vspace{0.2em}\noindent{\footnotesize\textcolor{gray}{%
      [topic: #1 \textbar{} level: #2 \textbar{} difficulty: #3 \textbar{} blueprint cell: #4]}}\par
  \fi
}

% Rationale, key only.
\newcommand{\rationale}[1]{\ifprintanswers\begin{solution}#1\end{solution}\fi}

% Keep each quiz item together on one page. A stem separated from its options,
% or options split across a page break, lets a student answer from a partial
% view of the alternatives -- which silently changes the item.
%
% \question cannot be wrapped: exam.cls (re)defines it inside the questions
% environment, so any outer redefinition is discarded. Items therefore call
% \qitem, which resolves \question at expansion time -- i.e. inside the
% environment, where it exists. \filbreak adapts to the item's actual height,
% which \Needspace's fixed line count cannot.
\newcommand{\qitem}{\filbreak\question}

\renewcommand{\choiceshook}{%
  \setlength\leftmargin{0.7em}%
  \setlength\itemsep{0.15em}%
  \setlength\topsep{0.3em}%
}
\renewcommand{\questionshook}{%
  \setlength\itemsep{0.9em}%
}

% Lowercase (a)/(b)/(c)/(d) option labels, matching how the answer key and the
% coding sheet refer to them.
\renewcommand{\thechoice}{\alph{choice}}
\renewcommand{\choicelabel}{(\thechoice)}

\checkboxchar{$\square$}
\checkedchar{$\blacksquare$}
\CorrectChoiceEmphasis{\bfseries}
\SolutionEmphasis{\small\itshape}
.
%   \surveywave  -> "Pre-Test" / "Post-Test"
%   \surveyform  -> "A" / "B"
% ---------------------------------------------------------------------------
\providecommand{\surveywave}{Survey}
\providecommand{\surveyform}{A}

\newcommand{\labtitle}{AI Engineering Lab}

% ---------------------------------------------------------------------------
% Running head / foot
% ---------------------------------------------------------------------------
\pagestyle{headandfoot}
\firstpageheader{}{}{}
\runningheader{\small\labtitle}{}{\small\surveywave{} --- Form \surveyform}
\runningheadrule
\firstpagefooter{}{\small Page \thepage\ of \numpages}{}
\runningfooter{}{\small Page \thepage\ of \numpages}{}

% No per-item point values: every quiz item is scored 1/0 externally
% (see docs/analysis_plan.md), so the printed form carries no point marks.
\nopointsinmargin

% ---------------------------------------------------------------------------
% Section headings for the questionnaire parts
% ---------------------------------------------------------------------------
\newcommand{\partheading}[2]{%
  \par\vspace{1.1em}%
  \noindent\begin{tabularx}{\linewidth}{@{}>{\bfseries}l X@{}}
    \large Part #1 & \large\bfseries #2
  \end{tabularx}%
  \par\vspace{-0.4em}\noindent\rule{\linewidth}{1pt}\par\vspace{0.5em}%
}

\newcommand{\blockheading}[1]{%
  \par\vspace{0.8em}\noindent\textbf{#1}%
  \par\vspace{-0.3em}\noindent\rule{\linewidth}{0.4pt}\par\vspace{0.3em}%
}

\newcommand{\instruction}[1]{\par\vspace{0.2em}\noindent\textit{\small #1}\par\vspace{0.5em}}

% A block heading *inside* the questions environment. Plain \blockheading cannot
% be used there: the questions environment is a list, and material placed in a
% list before or between items has to be escaped through \fullwidth.
\newcommand{\quizblock}[1]{\fullwidth{\blockheading{#1}}}

% ---------------------------------------------------------------------------
% Answer primitives
% ---------------------------------------------------------------------------
\newcommand{\bx}{$\square$}                        % empty checkbox
\newcommand{\bxl}[1]{\bx~#1}                       % checkbox with label
\newcommand{\writeline}[1][\linewidth]{\rule[-0.4em]{#1}{0.4pt}}

% Single-character write-in boxes, e.g. for the pseudonym code.
\newcommand{\charbox}{%
  \tikz[baseline=(c.base)]{\node[draw,minimum width=8mm,minimum height=9mm,inner sep=0pt] (c) {\phantom{X}};}%
}
\newcommand{\charboxes}[1]{%
  \foreach \i in {1,...,#1} {\charbox\hspace{1.2mm}}%
}

% ---------------------------------------------------------------------------
% Rating grids
% ---------------------------------------------------------------------------
% These use plain `tabular` with an explicitly computed p{} label column rather
% than `tabularx`. tabularx determines its body by scanning the input for a
% literal \end{tabularx}, which it cannot find when the \begin/\end pair is
% split across an environment definition -- so tabularx must not be wrapped in
% a \newenvironment. The label width is therefore derived from \linewidth minus
% the total width of the fixed check-box columns.

\setlength{\extrarowheight}{0pt}
\setlength{\emergencystretch}{2em}
\newlength{\gridlabel}
\newcommand{\ck}{\makebox[8mm]{$\square$}}          % one check-box column
\newcommand{\hd}[1]{\makebox[8mm]{\scriptsize #1}}  % matching header cell

% Forced-choice 1-5 rating scale. There is deliberately no "n/a" / "cannot judge"
% column: every statement here is about the respondent's own knowledge or
% experience, which they are always in a position to place on the scale. An
% escape column on such a grid is used mainly to skip, and each skip removes a
% repeated measure that cannot be recovered afterwards. A genuinely unanswerable
% row can still be left blank, which is recorded as missing.
%
% The legend is a mandatory argument because the grids do not all use the same
% scale: agreement and ability need different anchors, and a printed legend that
% does not match the question is worse than none.
% Usage: \begin{likert}[<column header>]{<legend>} \litem{stem} ... \end{likert}
\newcommand{\legendagree}{1 = strongly disagree \quad 2 = disagree \quad 3 = neutral \quad 4 = agree
  \quad 5 = strongly agree}
\newcommand{\legendability}{1 = no experience \quad 2 = beginner \quad 3 = competent
  \quad 4 = proficient \quad 5 = expert}
\newenvironment{likert}[2][Today, I would say \dots]{%
  \par\small
  % The legend is emitted after the table, where the environment's arguments are
  % no longer in scope, so it is stashed here.
  \def\likertlegend{#2}%
  \setlength{\tabcolsep}{2pt}%
  \setlength{\gridlabel}{\dimexpr\linewidth-56mm\relax}%
  \renewcommand{\arraystretch}{1.3}%
  \noindent
  \tabular{@{}p{\gridlabel} ccccc@{}}
    \toprule
    \textbf{#1} & \hd{1} & \hd{2} & \hd{3} & \hd{4} & \hd{5} \\
    \midrule
}{%
    \bottomrule
  \endtabular
  \par\vspace{0.35em}%
  {\scriptsize \likertlegend}%
  \par
}
\newcommand{\litem}[1]{#1 & \ck & \ck & \ck & \ck & \ck \\}

% Two-scale component grid: "can explain" and "can implement" side by side.
\newenvironment{dualgrid}[2]{%
  \par\small
  \setlength{\tabcolsep}{1.6pt}%
  \setlength{\gridlabel}{\dimexpr\linewidth-97mm\relax}%
  \renewcommand{\arraystretch}{1.3}%
  \noindent
  \tabular{@{}p{\gridlabel} @{\hspace{2mm}} ccccc @{\hspace{4mm}} ccccc@{}}
    \toprule
    & \multicolumn{5}{c}{\textbf{#1}} & \multicolumn{5}{c}{\textbf{#2}} \\
    \cmidrule(lr){2-6}\cmidrule(lr){7-11}
    \textbf{Component} & \hd{1} & \hd{2} & \hd{3} & \hd{4} & \hd{5}
                       & \hd{1} & \hd{2} & \hd{3} & \hd{4} & \hd{5} \\
    \midrule
}{%
    \bottomrule
  \endtabular
  \par\vspace{0.35em}%
  {\scriptsize 1 = not at all \quad 2 = a little \quad 3 = partly \quad 4 = mostly \quad 5 = fully}%
  \par
}
\newcommand{\ditem}[1]{#1 & \ck & \ck & \ck & \ck & \ck & \ck & \ck & \ck & \ck & \ck \\}

% Frequency grid: daily / weekly / monthly / tried once / never.
\newenvironment{freqgrid}[1]{%
  \par\small
  \setlength{\tabcolsep}{2pt}%
  \setlength{\gridlabel}{\dimexpr\linewidth-75mm\relax}%
  \renewcommand{\arraystretch}{1.3}%
  \noindent
  \tabular{@{}p{\gridlabel} ccccc@{}}
    \toprule
    \textbf{#1} & \makebox[12mm]{\scriptsize daily} & \makebox[12mm]{\scriptsize weekly}
                & \makebox[14mm]{\scriptsize monthly} & \makebox[16mm]{\scriptsize tried once}
                & \makebox[12mm]{\scriptsize never} \\
    \midrule
}{%
    \bottomrule
  \endtabular\par
}
\newcommand{\fitem}[1]{#1 & \makebox[12mm]{\bx} & \makebox[12mm]{\bx} & \makebox[14mm]{\bx}
                          & \makebox[16mm]{\bx} & \makebox[12mm]{\bx} \\}

% Experience grid: yes / attempted / no.
\newenvironment{expgrid}[1]{%
  \par\small
  \setlength{\tabcolsep}{2pt}%
  \setlength{\gridlabel}{\dimexpr\linewidth-54mm\relax}%
  \renewcommand{\arraystretch}{1.3}%
  \noindent
  \tabular{@{}p{\gridlabel} ccc@{}}
    \toprule
    \textbf{#1} & \makebox[12mm]{\scriptsize yes} & \makebox[22mm]{\scriptsize attempted}
                & \makebox[12mm]{\scriptsize no} \\
    \midrule
}{%
    \bottomrule
  \endtabular\par
}
\newcommand{\eitem}[1]{#1 & \makebox[12mm]{\bx} & \makebox[22mm]{\bx} & \makebox[12mm]{\bx} \\}

% ---------------------------------------------------------------------------
% Quiz items
% ---------------------------------------------------------------------------
% The "don't know" escape hatch. Placed after the choices so that the item
% itself remains a four-option single-best-answer item, while omission is
% still distinguishable from guessing in the data.
\newcommand{\dontknow}{%
  \par\vspace{-0.2em}\noindent\hspace{1.2em}{\small\bx~\textit{I do not know.}}\par
}

% Metadata shown only in the answer key.
\newcommand{\itemmeta}[4]{%
  \ifprintanswers
    \par\vspace{0.2em}\noindent{\footnotesize\textcolor{gray}{%
      [topic: #1 \textbar{} level: #2 \textbar{} difficulty: #3 \textbar{} blueprint cell: #4]}}\par
  \fi
}

% Rationale, key only.
\newcommand{\rationale}[1]{\ifprintanswers\begin{solution}#1\end{solution}\fi}

% Keep each quiz item together on one page. A stem separated from its options,
% or options split across a page break, lets a student answer from a partial
% view of the alternatives -- which silently changes the item.
%
% \question cannot be wrapped: exam.cls (re)defines it inside the questions
% environment, so any outer redefinition is discarded. Items therefore call
% \qitem, which resolves \question at expansion time -- i.e. inside the
% environment, where it exists. \filbreak adapts to the item's actual height,
% which \Needspace's fixed line count cannot.
\newcommand{\qitem}{\filbreak\question}

\renewcommand{\choiceshook}{%
  \setlength\leftmargin{0.7em}%
  \setlength\itemsep{0.15em}%
  \setlength\topsep{0.3em}%
}
\renewcommand{\questionshook}{%
  \setlength\itemsep{0.9em}%
}

% Lowercase (a)/(b)/(c)/(d) option labels, matching how the answer key and the
% coding sheet refer to them.
\renewcommand{\thechoice}{\alph{choice}}
\renewcommand{\choicelabel}{(\thechoice)}

\checkboxchar{$\square$}
\checkedchar{$\blacksquare$}
\CorrectChoiceEmphasis{\bfseries}
\SolutionEmphasis{\small\itshape}
.
%   \surveywave  -> "Pre-Test" / "Post-Test"
%   \surveyform  -> "A" / "B"
% ---------------------------------------------------------------------------
\providecommand{\surveywave}{Survey}
\providecommand{\surveyform}{A}

\newcommand{\labtitle}{AI Engineering Lab}

% ---------------------------------------------------------------------------
% Running head / foot
% ---------------------------------------------------------------------------
\pagestyle{headandfoot}
\firstpageheader{}{}{}
\runningheader{\small\labtitle}{}{\small\surveywave{} --- Form \surveyform}
\runningheadrule
\firstpagefooter{}{\small Page \thepage\ of \numpages}{}
\runningfooter{}{\small Page \thepage\ of \numpages}{}

% No per-item point values: every quiz item is scored 1/0 externally
% (see docs/analysis_plan.md), so the printed form carries no point marks.
\nopointsinmargin

% ---------------------------------------------------------------------------
% Section headings for the questionnaire parts
% ---------------------------------------------------------------------------
\newcommand{\partheading}[2]{%
  \par\vspace{1.1em}%
  \noindent\begin{tabularx}{\linewidth}{@{}>{\bfseries}l X@{}}
    \large Part #1 & \large\bfseries #2
  \end{tabularx}%
  \par\vspace{-0.4em}\noindent\rule{\linewidth}{1pt}\par\vspace{0.5em}%
}

\newcommand{\blockheading}[1]{%
  \par\vspace{0.8em}\noindent\textbf{#1}%
  \par\vspace{-0.3em}\noindent\rule{\linewidth}{0.4pt}\par\vspace{0.3em}%
}

\newcommand{\instruction}[1]{\par\vspace{0.2em}\noindent\textit{\small #1}\par\vspace{0.5em}}

% A block heading *inside* the questions environment. Plain \blockheading cannot
% be used there: the questions environment is a list, and material placed in a
% list before or between items has to be escaped through \fullwidth.
\newcommand{\quizblock}[1]{\fullwidth{\blockheading{#1}}}

% ---------------------------------------------------------------------------
% Answer primitives
% ---------------------------------------------------------------------------
\newcommand{\bx}{$\square$}                        % empty checkbox
\newcommand{\bxl}[1]{\bx~#1}                       % checkbox with label
\newcommand{\writeline}[1][\linewidth]{\rule[-0.4em]{#1}{0.4pt}}

% Single-character write-in boxes, e.g. for the pseudonym code.
\newcommand{\charbox}{%
  \tikz[baseline=(c.base)]{\node[draw,minimum width=8mm,minimum height=9mm,inner sep=0pt] (c) {\phantom{X}};}%
}
\newcommand{\charboxes}[1]{%
  \foreach \i in {1,...,#1} {\charbox\hspace{1.2mm}}%
}

% ---------------------------------------------------------------------------
% Rating grids
% ---------------------------------------------------------------------------
% These use plain `tabular` with an explicitly computed p{} label column rather
% than `tabularx`. tabularx determines its body by scanning the input for a
% literal \end{tabularx}, which it cannot find when the \begin/\end pair is
% split across an environment definition -- so tabularx must not be wrapped in
% a \newenvironment. The label width is therefore derived from \linewidth minus
% the total width of the fixed check-box columns.

\setlength{\extrarowheight}{0pt}
\setlength{\emergencystretch}{2em}
\newlength{\gridlabel}
\newcommand{\ck}{\makebox[8mm]{$\square$}}          % one check-box column
\newcommand{\hd}[1]{\makebox[8mm]{\scriptsize #1}}  % matching header cell

% Forced-choice 1-5 rating scale. There is deliberately no "n/a" / "cannot judge"
% column: every statement here is about the respondent's own knowledge or
% experience, which they are always in a position to place on the scale. An
% escape column on such a grid is used mainly to skip, and each skip removes a
% repeated measure that cannot be recovered afterwards. A genuinely unanswerable
% row can still be left blank, which is recorded as missing.
%
% The legend is a mandatory argument because the grids do not all use the same
% scale: agreement and ability need different anchors, and a printed legend that
% does not match the question is worse than none.
% Usage: \begin{likert}[<column header>]{<legend>} \litem{stem} ... \end{likert}
\newcommand{\legendagree}{1 = strongly disagree \quad 2 = disagree \quad 3 = neutral \quad 4 = agree
  \quad 5 = strongly agree}
\newcommand{\legendability}{1 = no experience \quad 2 = beginner \quad 3 = competent
  \quad 4 = proficient \quad 5 = expert}
\newenvironment{likert}[2][Today, I would say \dots]{%
  \par\small
  % The legend is emitted after the table, where the environment's arguments are
  % no longer in scope, so it is stashed here.
  \def\likertlegend{#2}%
  \setlength{\tabcolsep}{2pt}%
  \setlength{\gridlabel}{\dimexpr\linewidth-56mm\relax}%
  \renewcommand{\arraystretch}{1.3}%
  \noindent
  \tabular{@{}p{\gridlabel} ccccc@{}}
    \toprule
    \textbf{#1} & \hd{1} & \hd{2} & \hd{3} & \hd{4} & \hd{5} \\
    \midrule
}{%
    \bottomrule
  \endtabular
  \par\vspace{0.35em}%
  {\scriptsize \likertlegend}%
  \par
}
\newcommand{\litem}[1]{#1 & \ck & \ck & \ck & \ck & \ck \\}

% Two-scale component grid: "can explain" and "can implement" side by side.
\newenvironment{dualgrid}[2]{%
  \par\small
  \setlength{\tabcolsep}{1.6pt}%
  \setlength{\gridlabel}{\dimexpr\linewidth-97mm\relax}%
  \renewcommand{\arraystretch}{1.3}%
  \noindent
  \tabular{@{}p{\gridlabel} @{\hspace{2mm}} ccccc @{\hspace{4mm}} ccccc@{}}
    \toprule
    & \multicolumn{5}{c}{\textbf{#1}} & \multicolumn{5}{c}{\textbf{#2}} \\
    \cmidrule(lr){2-6}\cmidrule(lr){7-11}
    \textbf{Component} & \hd{1} & \hd{2} & \hd{3} & \hd{4} & \hd{5}
                       & \hd{1} & \hd{2} & \hd{3} & \hd{4} & \hd{5} \\
    \midrule
}{%
    \bottomrule
  \endtabular
  \par\vspace{0.35em}%
  {\scriptsize 1 = not at all \quad 2 = a little \quad 3 = partly \quad 4 = mostly \quad 5 = fully}%
  \par
}
\newcommand{\ditem}[1]{#1 & \ck & \ck & \ck & \ck & \ck & \ck & \ck & \ck & \ck & \ck \\}

% Frequency grid: daily / weekly / monthly / tried once / never.
\newenvironment{freqgrid}[1]{%
  \par\small
  \setlength{\tabcolsep}{2pt}%
  \setlength{\gridlabel}{\dimexpr\linewidth-75mm\relax}%
  \renewcommand{\arraystretch}{1.3}%
  \noindent
  \tabular{@{}p{\gridlabel} ccccc@{}}
    \toprule
    \textbf{#1} & \makebox[12mm]{\scriptsize daily} & \makebox[12mm]{\scriptsize weekly}
                & \makebox[14mm]{\scriptsize monthly} & \makebox[16mm]{\scriptsize tried once}
                & \makebox[12mm]{\scriptsize never} \\
    \midrule
}{%
    \bottomrule
  \endtabular\par
}
\newcommand{\fitem}[1]{#1 & \makebox[12mm]{\bx} & \makebox[12mm]{\bx} & \makebox[14mm]{\bx}
                          & \makebox[16mm]{\bx} & \makebox[12mm]{\bx} \\}

% Experience grid: yes / attempted / no.
\newenvironment{expgrid}[1]{%
  \par\small
  \setlength{\tabcolsep}{2pt}%
  \setlength{\gridlabel}{\dimexpr\linewidth-54mm\relax}%
  \renewcommand{\arraystretch}{1.3}%
  \noindent
  \tabular{@{}p{\gridlabel} ccc@{}}
    \toprule
    \textbf{#1} & \makebox[12mm]{\scriptsize yes} & \makebox[22mm]{\scriptsize attempted}
                & \makebox[12mm]{\scriptsize no} \\
    \midrule
}{%
    \bottomrule
  \endtabular\par
}
\newcommand{\eitem}[1]{#1 & \makebox[12mm]{\bx} & \makebox[22mm]{\bx} & \makebox[12mm]{\bx} \\}

% ---------------------------------------------------------------------------
% Quiz items
% ---------------------------------------------------------------------------
% The "don't know" escape hatch. Placed after the choices so that the item
% itself remains a four-option single-best-answer item, while omission is
% still distinguishable from guessing in the data.
\newcommand{\dontknow}{%
  \par\vspace{-0.2em}\noindent\hspace{1.2em}{\small\bx~\textit{I do not know.}}\par
}

% Metadata shown only in the answer key.
\newcommand{\itemmeta}[4]{%
  \ifprintanswers
    \par\vspace{0.2em}\noindent{\footnotesize\textcolor{gray}{%
      [topic: #1 \textbar{} level: #2 \textbar{} difficulty: #3 \textbar{} blueprint cell: #4]}}\par
  \fi
}

% Rationale, key only.
\newcommand{\rationale}[1]{\ifprintanswers\begin{solution}#1\end{solution}\fi}

% Keep each quiz item together on one page. A stem separated from its options,
% or options split across a page break, lets a student answer from a partial
% view of the alternatives -- which silently changes the item.
%
% \question cannot be wrapped: exam.cls (re)defines it inside the questions
% environment, so any outer redefinition is discarded. Items therefore call
% \qitem, which resolves \question at expansion time -- i.e. inside the
% environment, where it exists. \filbreak adapts to the item's actual height,
% which \Needspace's fixed line count cannot.
\newcommand{\qitem}{\filbreak\question}

\renewcommand{\choiceshook}{%
  \setlength\leftmargin{0.7em}%
  \setlength\itemsep{0.15em}%
  \setlength\topsep{0.3em}%
}
\renewcommand{\questionshook}{%
  \setlength\itemsep{0.9em}%
}

% Lowercase (a)/(b)/(c)/(d) option labels, matching how the answer key and the
% coding sheet refer to them.
\renewcommand{\thechoice}{\alph{choice}}
\renewcommand{\choicelabel}{(\thechoice)}

\checkboxchar{$\square$}
\checkedchar{$\blacksquare$}
\CorrectChoiceEmphasis{\bfseries}
\SolutionEmphasis{\small\itshape}
.
%   \surveywave  -> "Pre-Test" / "Post-Test"
%   \surveyform  -> "A" / "B"
% ---------------------------------------------------------------------------
\providecommand{\surveywave}{Survey}
\providecommand{\surveyform}{A}

\newcommand{\labtitle}{AI Engineering Lab}

% ---------------------------------------------------------------------------
% Running head / foot
% ---------------------------------------------------------------------------
\pagestyle{headandfoot}
\firstpageheader{}{}{}
\runningheader{\small\labtitle}{}{\small\surveywave{} --- Form \surveyform}
\runningheadrule
\firstpagefooter{}{\small Page \thepage\ of \numpages}{}
\runningfooter{}{\small Page \thepage\ of \numpages}{}

% No per-item point values: every quiz item is scored 1/0 externally
% (see docs/analysis_plan.md), so the printed form carries no point marks.
\nopointsinmargin

% ---------------------------------------------------------------------------
% Section headings for the questionnaire parts
% ---------------------------------------------------------------------------
\newcommand{\partheading}[2]{%
  \par\vspace{1.1em}%
  \noindent\begin{tabularx}{\linewidth}{@{}>{\bfseries}l X@{}}
    \large Part #1 & \large\bfseries #2
  \end{tabularx}%
  \par\vspace{-0.4em}\noindent\rule{\linewidth}{1pt}\par\vspace{0.5em}%
}

\newcommand{\blockheading}[1]{%
  \par\vspace{0.8em}\noindent\textbf{#1}%
  \par\vspace{-0.3em}\noindent\rule{\linewidth}{0.4pt}\par\vspace{0.3em}%
}

\newcommand{\instruction}[1]{\par\vspace{0.2em}\noindent\textit{\small #1}\par\vspace{0.5em}}

% A block heading *inside* the questions environment. Plain \blockheading cannot
% be used there: the questions environment is a list, and material placed in a
% list before or between items has to be escaped through \fullwidth.
\newcommand{\quizblock}[1]{\fullwidth{\blockheading{#1}}}

% ---------------------------------------------------------------------------
% Answer primitives
% ---------------------------------------------------------------------------
\newcommand{\bx}{$\square$}                        % empty checkbox
\newcommand{\bxl}[1]{\bx~#1}                       % checkbox with label
\newcommand{\writeline}[1][\linewidth]{\rule[-0.4em]{#1}{0.4pt}}

% Single-character write-in boxes, e.g. for the pseudonym code.
\newcommand{\charbox}{%
  \tikz[baseline=(c.base)]{\node[draw,minimum width=8mm,minimum height=9mm,inner sep=0pt] (c) {\phantom{X}};}%
}
\newcommand{\charboxes}[1]{%
  \foreach \i in {1,...,#1} {\charbox\hspace{1.2mm}}%
}

% ---------------------------------------------------------------------------
% Rating grids
% ---------------------------------------------------------------------------
% These use plain `tabular` with an explicitly computed p{} label column rather
% than `tabularx`. tabularx determines its body by scanning the input for a
% literal \end{tabularx}, which it cannot find when the \begin/\end pair is
% split across an environment definition -- so tabularx must not be wrapped in
% a \newenvironment. The label width is therefore derived from \linewidth minus
% the total width of the fixed check-box columns.

\setlength{\extrarowheight}{0pt}
\setlength{\emergencystretch}{2em}
\newlength{\gridlabel}
\newcommand{\ck}{\makebox[8mm]{$\square$}}          % one check-box column
\newcommand{\hd}[1]{\makebox[8mm]{\scriptsize #1}}  % matching header cell

% Forced-choice 1-5 rating scale. There is deliberately no "n/a" / "cannot judge"
% column: every statement here is about the respondent's own knowledge or
% experience, which they are always in a position to place on the scale. An
% escape column on such a grid is used mainly to skip, and each skip removes a
% repeated measure that cannot be recovered afterwards. A genuinely unanswerable
% row can still be left blank, which is recorded as missing.
%
% The legend is a mandatory argument because the grids do not all use the same
% scale: agreement and ability need different anchors, and a printed legend that
% does not match the question is worse than none.
% Usage: \begin{likert}[<column header>]{<legend>} \litem{stem} ... \end{likert}
\newcommand{\legendagree}{1 = strongly disagree \quad 2 = disagree \quad 3 = neutral \quad 4 = agree
  \quad 5 = strongly agree}
\newcommand{\legendability}{1 = no experience \quad 2 = beginner \quad 3 = competent
  \quad 4 = proficient \quad 5 = expert}
\newenvironment{likert}[2][Today, I would say \dots]{%
  \par\small
  % The legend is emitted after the table, where the environment's arguments are
  % no longer in scope, so it is stashed here.
  \def\likertlegend{#2}%
  \setlength{\tabcolsep}{2pt}%
  \setlength{\gridlabel}{\dimexpr\linewidth-56mm\relax}%
  \renewcommand{\arraystretch}{1.3}%
  \noindent
  \tabular{@{}p{\gridlabel} ccccc@{}}
    \toprule
    \textbf{#1} & \hd{1} & \hd{2} & \hd{3} & \hd{4} & \hd{5} \\
    \midrule
}{%
    \bottomrule
  \endtabular
  \par\vspace{0.35em}%
  {\scriptsize \likertlegend}%
  \par
}
\newcommand{\litem}[1]{#1 & \ck & \ck & \ck & \ck & \ck \\}

% Two-scale component grid: "can explain" and "can implement" side by side.
\newenvironment{dualgrid}[2]{%
  \par\small
  \setlength{\tabcolsep}{1.6pt}%
  \setlength{\gridlabel}{\dimexpr\linewidth-97mm\relax}%
  \renewcommand{\arraystretch}{1.3}%
  \noindent
  \tabular{@{}p{\gridlabel} @{\hspace{2mm}} ccccc @{\hspace{4mm}} ccccc@{}}
    \toprule
    & \multicolumn{5}{c}{\textbf{#1}} & \multicolumn{5}{c}{\textbf{#2}} \\
    \cmidrule(lr){2-6}\cmidrule(lr){7-11}
    \textbf{Component} & \hd{1} & \hd{2} & \hd{3} & \hd{4} & \hd{5}
                       & \hd{1} & \hd{2} & \hd{3} & \hd{4} & \hd{5} \\
    \midrule
}{%
    \bottomrule
  \endtabular
  \par\vspace{0.35em}%
  {\scriptsize 1 = not at all \quad 2 = a little \quad 3 = partly \quad 4 = mostly \quad 5 = fully}%
  \par
}
\newcommand{\ditem}[1]{#1 & \ck & \ck & \ck & \ck & \ck & \ck & \ck & \ck & \ck & \ck \\}

% Frequency grid: daily / weekly / monthly / tried once / never.
\newenvironment{freqgrid}[1]{%
  \par\small
  \setlength{\tabcolsep}{2pt}%
  \setlength{\gridlabel}{\dimexpr\linewidth-75mm\relax}%
  \renewcommand{\arraystretch}{1.3}%
  \noindent
  \tabular{@{}p{\gridlabel} ccccc@{}}
    \toprule
    \textbf{#1} & \makebox[12mm]{\scriptsize daily} & \makebox[12mm]{\scriptsize weekly}
                & \makebox[14mm]{\scriptsize monthly} & \makebox[16mm]{\scriptsize tried once}
                & \makebox[12mm]{\scriptsize never} \\
    \midrule
}{%
    \bottomrule
  \endtabular\par
}
\newcommand{\fitem}[1]{#1 & \makebox[12mm]{\bx} & \makebox[12mm]{\bx} & \makebox[14mm]{\bx}
                          & \makebox[16mm]{\bx} & \makebox[12mm]{\bx} \\}

% Experience grid: yes / attempted / no.
\newenvironment{expgrid}[1]{%
  \par\small
  \setlength{\tabcolsep}{2pt}%
  \setlength{\gridlabel}{\dimexpr\linewidth-54mm\relax}%
  \renewcommand{\arraystretch}{1.3}%
  \noindent
  \tabular{@{}p{\gridlabel} ccc@{}}
    \toprule
    \textbf{#1} & \makebox[12mm]{\scriptsize yes} & \makebox[22mm]{\scriptsize attempted}
                & \makebox[12mm]{\scriptsize no} \\
    \midrule
}{%
    \bottomrule
  \endtabular\par
}
\newcommand{\eitem}[1]{#1 & \makebox[12mm]{\bx} & \makebox[22mm]{\bx} & \makebox[12mm]{\bx} \\}

% ---------------------------------------------------------------------------
% Quiz items
% ---------------------------------------------------------------------------
% The "don't know" escape hatch. Placed after the choices so that the item
% itself remains a four-option single-best-answer item, while omission is
% still distinguishable from guessing in the data.
\newcommand{\dontknow}{%
  \par\vspace{-0.2em}\noindent\hspace{1.2em}{\small\bx~\textit{I do not know.}}\par
}

% Metadata shown only in the answer key.
\newcommand{\itemmeta}[4]{%
  \ifprintanswers
    \par\vspace{0.2em}\noindent{\footnotesize\textcolor{gray}{%
      [topic: #1 \textbar{} level: #2 \textbar{} difficulty: #3 \textbar{} blueprint cell: #4]}}\par
  \fi
}

% Rationale, key only.
\newcommand{\rationale}[1]{\ifprintanswers\begin{solution}#1\end{solution}\fi}

% Keep each quiz item together on one page. A stem separated from its options,
% or options split across a page break, lets a student answer from a partial
% view of the alternatives -- which silently changes the item.
%
% \question cannot be wrapped: exam.cls (re)defines it inside the questions
% environment, so any outer redefinition is discarded. Items therefore call
% \qitem, which resolves \question at expansion time -- i.e. inside the
% environment, where it exists. \filbreak adapts to the item's actual height,
% which \Needspace's fixed line count cannot.
\newcommand{\qitem}{\filbreak\question}

\renewcommand{\choiceshook}{%
  \setlength\leftmargin{0.7em}%
  \setlength\itemsep{0.15em}%
  \setlength\topsep{0.3em}%
}
\renewcommand{\questionshook}{%
  \setlength\itemsep{0.9em}%
}

% Lowercase (a)/(b)/(c)/(d) option labels, matching how the answer key and the
% coding sheet refer to them.
\renewcommand{\thechoice}{\alph{choice}}
\renewcommand{\choicelabel}{(\thechoice)}

\checkboxchar{$\square$}
\checkedchar{$\blacksquare$}
\CorrectChoiceEmphasis{\bfseries}
\SolutionEmphasis{\small\itshape}
